% ============================================================
% Cole cada \section abaixo no lugar certo do samplepaper.tex.
% Antes de \begin{document}, garanta que o preamble tem:
%   \usepackage{graphicx}
%   \usepackage{booktabs}
% Suba as 5 imagens pro Overleaf (aba de arquivos, "Upload"):
%   streamlit_app.png, software_architecture.png,
%   architecture_before.png, architecture_after.png,
%   (estrutura_pastas.png e git_gitflow.png vão no OUTRO slide, não aqui)
% ============================================================

\section{Introdução}

Este relatório descreve a implementação de uma rede neural multicamadas
(MLP -- \textit{Multi-Layer Perceptron}) construída do zero, representada
explicitamente como um grafo de neurônios e conexões, sem uso de bibliotecas
de \textit{autograd} (PyTorch, TensorFlow). A motivação da prática é
consolidar, na prática, os mecanismos de \textit{feed forward} e
\textit{back propagation} que sustentam o treinamento de redes neurais,
tornando explícitos os passos que normalmente ficam ocultos por trás de
chamadas de alto nível em frameworks de deep learning.

A especificação da prática pedia: (i) codificar uma MLP a partir da
estrutura de um grafo; (ii) implementar manualmente feed forward e back
propagation; (iii) treinar e avaliar a rede na base \textit{Heart Disease}
(divisão 80\% treino / 20\% teste); e (iv) demonstrar/visualizar o processo
de treinamento e inferência, com o grafo e os pesos sendo atualizados.
Este relatório documenta como cada um desses pontos foi resolvido, os
resultados obtidos e as variações de hiperparâmetros testadas.

\section{Método}

\subsection{Etapas, bibliotecas e uso de IA}

O projeto foi construído em etapas: (1) leitura da especificação da
prática (estrutura de grafo, feed forward/back propagation manuais,
dataset e split 80/20); (2) implementação dos módulos de base -- grafo,
ativações, forward/backward; (3) validação da implementação por
\textit{gradient checking}; (4) obtenção e preparo do dataset
\textit{Heart Disease}; (5) busca manual de hiperparâmetros (arquitetura,
ativação, \textit{learning rate}, número de épocas -- detalhada na seção
de Resultados); (6) construção do harness interativo em Streamlit;
(7) geração dos diagramas e redação deste relatório.

As bibliotecas usadas foram deliberadamente mínimas: \texttt{networkx}
para representar o grafo da rede, \texttt{matplotlib} para desenhá-lo
(estático e animado) e para os diagramas deste relatório,
\texttt{streamlit} para o harness web, e \texttt{unittest} (biblioteca
padrão) para os testes de \textit{gradient checking}. O núcleo da MLP
(\texttt{graph.py}, \texttt{activations.py}, \texttt{mlp.py}) não usa
NumPy, PyTorch, TensorFlow nem \textit{pandas} -- decisão deliberada, já
que o objetivo da prática é implementar feed forward e back propagation
de forma explícita, nó a nó e aresta a aresta, sem escondê-los por trás
de operações vetoriais ou de \textit{autograd}.

O desenvolvimento foi assistido pelo \textbf{Claude Code} (Anthropic), um
assistente de codificação por linha de comando. A IA foi usada para:
gerar a implementação inicial de cada módulo a partir da especificação da
prática, escrever os testes de \textit{gradient checking}, construir o
harness em Streamlit, gerar os diagramas de arquitetura (Figuras~\ref{fig:arquitetura}
e~\ref{fig:antes_depois}) e redigir o rascunho deste relatório a partir
dos resultados obtidos. As decisões de design (representar a rede como
grafo explícito, dividir o código em módulos pequenos), a escolha e
comparação dos hiperparâmetros (seção de Variações testadas) e a
validação dos resultados foram conduzidas e revisadas pelo autor a cada
etapa -- a IA funcionou como ferramenta de implementação e documentação,
não como autora das decisões experimentais.

\subsection{Organização da solução}

O projeto está organizado em módulos pequenos e independentes, cada um com
uma responsabilidade única: carregamento e pré-processamento dos dados,
construção do grafo, execução de feed forward/back propagation sobre esse
grafo, visualização, e dois pontos de entrada (um script de treino via
linha de comando e um harness interativo). A Figura~\ref{fig:arquitetura}
mostra os módulos e como eles se conectam.

\begin{figure}[htbp]
    \centering
    \includegraphics[width=0.85\textwidth]{software_architecture.png}
    \caption{Módulos do projeto e suas dependências. Os módulos de base
    (\texttt{data.py}, \texttt{graph.py}, \texttt{activations.py}) alimentam
    o motor de treino (\texttt{mlp.py}) e a visualização
    (\texttt{visualize.py}), que por sua vez são usados pelos dois pontos de
    entrada: o script de treino em lote (\texttt{train.py}) e o harness
    interativo (\texttt{app\_streamlit.py}).}
    \label{fig:arquitetura}
\end{figure}

A rede é representada como um \texttt{DiGraph} (grafo direcionado) da
biblioteca \texttt{networkx}: cada neurônio é um nó, com atributos de soma
ponderada (\textit{z}), ativação (\textit{value}), erro retropropagado
(\textit{delta}) e \textit{bias}; cada conexão entre dois neurônios de
camadas adjacentes é uma aresta, com um peso. O \textit{feed forward}
percorre os nós camada a camada, somando, para cada nó, o valor de cada nó
de origem multiplicado pelo peso da aresta correspondente, e aplicando a
função de ativação. O \textit{back propagation} percorre o grafo no sentido
inverso: o erro do neurônio de saída (diferença entre a previsão e o rótulo
real) é distribuído para trás, aresta por aresta, proporcionalmente ao peso
de cada conexão e à derivada da função de ativação em cada nó -- produzindo
o gradiente de cada peso e de cada \textit{bias}, usados para atualizar os
parâmetros por gradiente descendente.

A corretude da implementação do \textit{back propagation} foi verificada
por \textit{gradient checking}: o gradiente analítico calculado pelo
backward é comparado ao gradiente numérico (diferenças finitas), obtido
perturbando cada peso em $\pm\epsilon$ e observando a variação da perda.
Os dois métodos convergem até a quarta casa decimal, para as ativações
ReLU e tanh testadas.

\subsection{Harness de visualização}

O harness de referência da disciplina (\texttt{lsfcin/mlp}) é um aplicativo
desktop em PyQt6. Para este projeto, foi construído um harness equivalente
como aplicação web, usando \textit{Streamlit}, reaproveitando exatamente os
mesmos módulos de grafo e treino (nenhuma lógica foi duplicada). O harness
permite: carregar o dataset, ajustar arquitetura/ativação/taxa de
aprendizagem, executar um passo de treino ou uma época completa, rodar a
inferência no conjunto de teste, e visualizar o grafo (nós e arestas
coloridas por sinal e espessura por magnitude do peso) e a curva de perda
em tempo real. A Figura~\ref{fig:harness} mostra o harness em execução,
após carregar o dataset e treinar algumas épocas.

\begin{figure}[htbp]
    \centering
    \includegraphics[width=0.95\textwidth]{streamlit_app.png}
    \caption{Harness interativo (Streamlit) em execução: painel de
    controle à esquerda, grafo da rede com os pesos ao centro, e curva de
    perda (loss) por época à direita.}
    \label{fig:harness}
\end{figure}

\subsection{Dados}

A base utilizada foi a \textit{Heart Disease} (13 atributos clínicos --
idade, sexo, tipo de dor no peito, pressão em repouso, colesterol, entre
outros -- e um rótulo binário indicando presença de doença cardíaca), com
303 amostras. Os dados foram divididos em 80\% treino (242 amostras) e 20\%
teste (60 amostras), com separação fixa por semente aleatória. As features
foram padronizadas (\textit{z-score}: subtração da média e divisão pelo
desvio padrão), usando estatísticas calculadas apenas no conjunto de
treino, para evitar vazamento de informação do teste.

\section{Resultados e Discussão}

\subsection{Configuração final e métricas}

A configuração que obteve o melhor resultado no conjunto de teste foi:
arquitetura 13$\,|\,$8$\,|\,$5$\,|\,$1 (13 entradas, duas camadas escondidas
de 8 e 5 neurônios, 1 saída), ativação ReLU nas camadas escondidas e
sigmoid na saída, taxa de aprendizagem 0.005, treinada por 60 épocas com
gradiente descendente estocástico (atualização a cada amostra). A
Tabela~\ref{tab:metricas} reporta as métricas no conjunto de teste.

\begin{table}[htbp]
    \centering
    \caption{Métricas no conjunto de teste (60 amostras).}
    \label{tab:metricas}
    \begin{tabular}{lc}
        \toprule
        Métrica & Valor \\
        \midrule
        Acurácia  & 0.833 \\
        Precisão  & 0.757 \\
        Revocação (\textit{recall}) & 0.966 \\
        F1-score  & 0.848 \\
        \midrule
        Verdadeiros positivos & 28 \\
        Falsos positivos      & 9  \\
        Verdadeiros negativos & 22 \\
        Falsos negativos      & 1  \\
        \bottomrule
    \end{tabular}
\end{table}

A revocação alta (0.966) indica que o modelo raramente deixa passar um
paciente com doença cardíaca (apenas 1 falso negativo em 29 casos
positivos) -- um comportamento desejável nesse domínio, em que o custo de
um falso negativo (não detectar a doença) tende a ser maior que o de um
falso positivo. Em contrapartida, a precisão mais baixa (0.757) mostra que
o modelo superestima casos positivos com alguma frequência (9 falsos
positivos).

\subsection{Arquitetura antes e depois do treinamento}

A Figura~\ref{fig:antes_depois} mostra o grafo da rede em dois momentos: os
pesos logo após a inicialização (Xavier/Glorot) e os pesos após as 60
épocas de treinamento. É possível notar que a distribuição de pesos muda
de forma visível, sobretudo nas conexões entre a última camada escondida e
o neurônio de saída, que passam a ter magnitudes bem mais desiguais entre
si -- indício de que a rede aprendeu a dar peso diferente a diferentes
combinações de atributos.

\begin{figure}[htbp]
    \centering
    \includegraphics[width=0.47\textwidth]{architecture_before.png}\hfill
    \includegraphics[width=0.47\textwidth]{architecture_after.png}
    \caption{Grafo da rede antes do treinamento (pesos inicializados por
    Xavier/Glorot, à esquerda) e depois de 60 épocas de treinamento
    (à direita). Azul indica peso positivo, vermelho indica peso negativo;
    a espessura da aresta é proporcional à magnitude do peso.}
    \label{fig:antes_depois}
\end{figure}

\subsection{Variações testadas}

Diversas combinações de arquitetura, ativação, taxa de aprendizagem e
número de épocas foram testadas manualmente antes de se chegar à
configuração final. A Tabela~\ref{tab:variacoes} resume os principais
experimentos.

\begin{table}[htbp]
    \centering
    \caption{Variações de hiperparâmetros testadas (acurácia em treino e teste).}
    \label{tab:variacoes}
    \begin{tabular}{lccccc}
        \toprule
        Arquitetura & Ativação & LR & Épocas & Acc. treino & Acc. teste \\
        \midrule
        13\,|\,8\,|\,5\,|\,1  & relu & 0.01  & 150 & 1.000 & 0.733 \\
        13\,|\,8\,|\,5\,|\,1  & relu & 0.01  & 60  & 0.967 & 0.767 \\
        13\,|\,8\,|\,5\,|\,1  & tanh & 0.01  & 40  & 0.942 & 0.800 \\
        13\,|\,6\,|\,4\,|\,1  & relu & 0.01  & 50  & 0.938 & 0.683 \\
        13\,|\,8\,|\,5\,|\,1  & relu & 0.005 & 80  & 0.951 & 0.817 \\
        13\,|\,10\,|\,6\,|\,1 & relu & 0.005 & 80  & 0.955 & 0.700 \\
        13\,|\,8\,|\,5\,|\,1  & relu & 0.005 & 100 & 0.971 & 0.783 \\
        \textbf{13\,|\,8\,|\,5\,|\,1}  & \textbf{relu} & \textbf{0.005} & \textbf{60}  & \textbf{0.930} & \textbf{0.833} \\
        \bottomrule
    \end{tabular}
\end{table}

O padrão mais consistente observado foi \textbf{overfitting} em redes
maiores ou treinadas por mais épocas: a primeira linha da
Tabela~\ref{tab:variacoes} (150 épocas) atinge 100\% de acurácia no
treino, mas cai para 73,3\% no teste -- a rede passou a memorizar os
242 pacientes de treino em vez de aprender um padrão que generaliza para
pacientes novos. O mesmo efeito, em menor grau, aparece ao aumentar a
arquitetura (13\,|\,10\,|\,6\,|\,1): mais parâmetros, mais capacidade de
memorização, pior generalização nesse dataset pequeno. A configuração
final (8\,|\,5, 60 épocas, \textit{learning rate} 0.005) foi a que
equilibrou melhor esse compromisso entre ajuste ao treino e generalização.
A padronização das features (\textit{z-score}) também se mostrou
necessária: sem ela, atributos de escala maior (por exemplo, colesterol,
que varia de 126 a 564) dominam a soma ponderada e tornam o treinamento
instável, independentemente da arquitetura escolhida.
